\chapter{Graphs through Their Eigenvalues}\label{linear:ch:graphs}

\begin{center}\itshape
In this paper, we show that every $(2^{n-1}+1)$-vertex induced subgraph of the $n$-dimensional hypercube has maximum degree at
least $\sqrt n$.
\upshape(Hao Huang, opening of the abstract of his 2019 paper in the \emph{Annals of Mathematics})
\end{center}

At a party, every pair of guests has exactly one common friend. Friendship goes both ways, and nobody counts as their own friend. Must somebody know everyone? Try drawing the friendships. Three guests can form a triangle; five can form two triangles sharing a centre. The friendship theorem says that every finite example has this shape: triangles joined at a single guest.

Erd\H{o}s, R\'enyi and S\'os proved the theorem in 1966. Its spectral proof turns the friendship rule into a matrix identity. A trace then forces a rational square root to be an integer. Aigner and Ziegler call the central guest the politician in \emph{Proofs from THE BOOK}.

A much larger matrix settled another problem in 2019. Huang assigned signs to the edges of the $n$-cube so that the resulting matrix squares to $nI$. Interlacing then forces a vertex of high degree in any induced subgraph containing more than half the cube. His paper connected this bound to the sensitivity conjecture for Boolean functions. The matrix argument itself needs only the spectral tools already proved here.

A graph's adjacency matrix is real symmetric. We may therefore compare eigenvalues, count dimensions of eigenspaces and test Rayleigh quotients. We also use the non-negative-matrix theory from Volume~3, Appendix~A: Perron's theorem (Theorem~A.1), Frobenius's theorem (Theorem~A.6), and the row-sum bound (Exercise~A.2). The arithmetic restrictions they reveal will rule out some graphs altogether.

\portrait{frobenius}{Ferdinand Georg Frobenius (1849--1917). Photograph from the album for Weierstrass, between 1870 and 1885. Public domain, via Wikimedia Commons.}

Other graph questions need different pieces of the spectrum. Chapter~\ref{ch:cheeger} relates a spectral gap to cuts and walks; Chapter~\ref{ch:expanders} studies families with a uniform gap. The matrix-tree theorem, which counts spanning trees by a determinant, belongs to Volume~12, Appendix~A. Kirchhoff's connection between electrical networks and trees is another early example of graphs yielding to linear algebra.

\portrait{kirchhoff}{Gustav Robert Kirchhoff (1824--1887). Nineteenth-century photograph, artist unknown; Smithsonian Libraries. Public domain, via Wikimedia Commons.}

\begin{think}
Draw friendship examples on three, five and seven vertices. After drawing one triangle, how can you add two guests without giving any pair a second common friend? What obstruction would you expect if all guests had the same number of friends and none knew everybody?
\end{think}

\section{Adjacency matrices and walks}

Our graphs are finite, simple and undirected, with a nonempty vertex set $V$ of size $n$. An edge is an unordered pair of distinct vertices; write $i\sim j$ when it is present. The \emph{adjacency matrix} has entries $a_{ij}=1$ for edges and $0$ otherwise. Symmetry gives real eigenvalues $\lambda_1\ge\cdots\ge\lambda_n$. Its zero diagonal gives $\sum_i\lambda_i=\operatorname{tr}A=0$. Changing vertex names conjugates $A$ by a permutation matrix, so its spectrum is a property of the graph.

Matrix multiplication counts walks: $(A^k)_{ij}$ is the number of length-$k$ walks from $i$ to $j$, as proved in Volume~3, Appendix~A. A closed walk of length two chooses a neighbour and returns. A closed walk of length three travels around a triangle, choosing one of three starting vertices and either direction. Thus
\[
\sum_i\lambda_i^2=2|E|,\qquad \sum_i\lambda_i^3=6t(G),
\]
where $t(G)$ counts triangles. The spectrum already determines the numbers of vertices, edges and triangles.

If every degree is $k$, call the graph \emph{$k$-regular}. Then $A\mathbf1=k\mathbf1$, and the row-sum bound from Volume~3, Exercise~A.2, says that $k$ is the largest eigenvalue. If the graph is connected, Frobenius's theorem (Volume~3, Theorem~A.6) makes that eigenvalue simple. Conversely, each connected component of a regular graph supplies its own eigenvector $\mathbf1_C$ with eigenvalue $k$, so simplicity forces connectedness. The other eigenvectors may be chosen in $\mathbf1^\perp$. The all-ones matrix $J$ vanishes there, a useful simplification in the coming identities.

\begin{example}[basic spectra]\label{linear:exa:spectra}
Four families admit explicit eigenbases.
\begin{enumerate}[label=(\alph*)]
\item For the complete graph $K_n$, $A=J-I$. The constant direction has eigenvalue $n-1$, and its orthogonal complement has eigenvalue $-1$, with multiplicity $n-1$.
\item For a cycle $C_n$ ($n\ge3$), write $A=S+S^{-1}$ using cyclic shift. With $\omega=e^{2\pi i/n}$, the vectors $(\omega^{jm})_m$ give eigenvalues $2\cos(2\pi j/n)$, for $j=0,\ldots,n-1$.
\item A path $P_n$ has eigenvalues $2\cos(j\pi/(n+1))$, for $j=1,\ldots,n$; the sine eigenvectors appear in Exercise~\ref{linear:ex:pathgraph}.
\item The cube $Q_n$ joins binary words differing in one coordinate. Its eigenvalues are $n-2j$, with multiplicities $\binom nj$, for $j=0,\ldots,n$. Exercise~\ref{linear:ex:hypercube} constructs the eigenvectors directly.
\end{enumerate}
\end{example}

\section{The friendship theorem}

\begin{theorem}[Erd\H{o}s--R\'enyi--S\'os, 1966]\label{linear:thm:friendship}
Suppose every two distinct vertices of a finite nonempty graph have exactly one common neighbour. Then the graph is a \emph{windmill}: it has $2m+1$ vertices and consists of $m$ triangles sharing one central vertex. The case $m=0$ is a single vertex.
\end{theorem}

\begin{figure}[ht]
\centering
\fig{linear-ch08-windmill}
\caption{The windmill with four blades: every two vertices have exactly one common neighbour, and the centre is the
politician.}
\label{linear:fig:windmill}
\end{figure}

\begin{proof}
A single vertex already has the required description. Assume there are at least two vertices.

\emph{Compare non-neighbours.} If $u\not\sim v$, assign to each neighbour $x$ of $u$ the unique common neighbour $f(x)$ of $x$ and $v$. This lies in the neighbourhood of $v$. If distinct $x,x'$ had the same image $y$, then $u$ and $y$ would have common neighbours $x,x'$. Here $y\ne u$, since $y\sim v$ while $u\not\sim v$. Thus the assignment is injective. Reversing $u,v$ shows that they have equal degrees.

\emph{Assume there is no central vertex.} In the complement graph, degrees in the original graph are constant along edges, and hence along connected components. If the complement were disconnected, none of its components could be a singleton: that vertex would be adjacent to everyone in the original graph. Choose two components, each with at least two vertices. Two vertices of one would share every vertex of the other as a neighbour in the original graph, contrary to the hypothesis. Thus the complement is connected and $G$ is $k$-regular. Every vertex has a neighbour, and the ends of an edge must have a common neighbour, so $k\ge2$.

\emph{Use the trace.} Off the diagonal, the entries of $A^2$ are all $1$; on the diagonal they are $k$. Hence
\[
A^2=(k-1)I+J.
\]
Applying this identity to $\mathbf1$ gives $k^2=k-1+n$, or $n-1=k(k-1)$. On $\mathbf1^\perp$, the eigenvalues of $A$ are $\sqrt{k-1}$ and $-\sqrt{k-1}$. Write their multiplicities as $r,s$. The trace gives
\[
k+(r-s)\sqrt{k-1}=0.
\]
Since $k>0$, $r\ne s$, so the square root is rational. A rational square root of an integer is an integer: reducing a fraction to lowest terms and squaring shows its denominator must be $1$. Put $h=\sqrt{k-1}\ge1$. The trace identity becomes $h(s-r)=h^2+1$, so $h$ divides $1$. It follows that $h=1$, $k=2$ and $n=3$. But this graph is a triangle, with a central vertex after all. The assumption was impossible.

\emph{Pair the remaining vertices.} Choose a vertex $w$ adjacent to everyone else. For each $u\ne w$, the unique common neighbour of $u,w$ is a vertex $u'$ adjacent to $u$. There can be no other neighbour of $u$ besides $w,u'$, since it would be a second common neighbour of $u,w$. Thus all vertices other than $w$ split into pairs, each completing a triangle with $w$.
\end{proof}

The last pairing argument did not use finiteness. The trace did. Exercise~\ref{linear:ex:infinitefriends} constructs an infinite friendship graph with no central vertex and shows exactly where this finite proof stops applying.

\section{Strongly regular graphs and Moore graphs}

\begin{definition}\label{linear:def:srg}
A $k$-regular graph with $n$ vertices is \emph{strongly regular with parameters $(n,k,\lambda,\mu)$} when an adjacent pair has exactly $\lambda$ common neighbours and a distinct non-adjacent pair has exactly $\mu$ common neighbours.
\end{definition}

The definition allows complete graphs, for which the non-adjacent-pair requirement is vacuous. We will separate that case when computing spectra. For noncomplete examples, the pentagon has parameters $(5,2,0,1)$ and the Petersen graph has $(10,3,0,1)$. Counting common neighbours in each matrix position gives
\begin{equation}\label{linear:eq:srg}
A^2=kI+\lambda A+\mu(J-I-A).
\end{equation}

\begin{figure}[ht]
\centering
\fig{linear-ch08-petersen}
\caption{The Petersen graph, strongly regular with parameters $(10,3,0,1)$: no triangles, and two non-adjacent vertices have
exactly one common neighbour. Its spectrum is $3$, $1$ (five times), $-2$ (four times). Its vertices can also be read as the
$2$-element subsets of $\{1,\dots,5\}$, adjacent when disjoint.}
\label{linear:fig:petersen}
\end{figure}

\begin{theorem}[eigenvalues of strongly regular graphs]\label{linear:thm:srg}
Let $G$ be connected and strongly regular with parameters $(n,k,\lambda,\mu)$ and $\mu>0$. If $G$ is complete, its spectrum is $k=n-1$ once and $-1$ with multiplicity $n-1$. Otherwise its spectrum consists of the simple eigenvalue $k$ and the roots $\theta>\tau$ of
\[
x^2-(\lambda-\mu)x-(k-\mu)=0.
\]
The respective root multiplicities are
\[
m_\theta,m_\tau=\frac12\left((n-1)\mp
\frac{2k+(n-1)(\lambda-\mu)}{\sqrt{(\lambda-\mu)^2+4(k-\mu)}}\right).
\]
In particular these expressions must be non-negative integers.
\end{theorem}
\begin{proof}
The complete case was computed in Example~\ref{linear:exa:spectra}. Otherwise connectedness makes $k$ simple, and every other eigenvector lies in $\mathbf1^\perp$. There $J=0$, so~\eqref{linear:eq:srg} forces the displayed quadratic. Nonadjacent vertices exist, so $\mu\le k$. Also $k>0$ and $\lambda\le k-1$. Consequently $\Delta=(\lambda-\mu)^2+4(k-\mu)>0$: if both terms vanished, we would have $\lambda=\mu=k$, a contradiction. The roots are distinct.

Their multiplicities satisfy $m_\theta+m_\tau=n-1$ and $k+m_\theta\theta+m_\tau\tau=0$. Substitute $\theta+\tau=\lambda-\mu$ and $\theta-\tau=\sqrt\Delta$ into the latter equation to obtain
\[
0=k+\frac{(n-1)(\lambda-\mu)}2+
\frac{(m_\theta-m_\tau)\sqrt\Delta}2.
\]
Solving the sum and difference equations gives the formulas. Multiplicities count eigenspace dimensions, so they must be non-negative integers.
\end{proof}

Starting at a vertex of degree $k$, we can reach at most $k$ vertices in one step and $k(k-1)$ more in two. A graph of diameter $2$ and maximum degree $k$ therefore has at most $k^2+1$ vertices. A regular graph attaining this bound is a \emph{Moore graph of diameter $2$}. Equality leaves no room for triangles, four-cycles or repeated second neighbours; Exercise~\ref{linear:ex:moorebound} proves these assertions. Its strongly regular parameters are $(k^2+1,k,0,1)$.

\begin{theorem}[Hoffman--Singleton, 1960]\label{linear:thm:moore}
A Moore graph of diameter $2$ can have only one of the degrees $2,3,7,57$.
\end{theorem}
\begin{proof}
Its parameters give $\Delta=4k-3$ and $m_\theta+m_\tau=k^2$. The trace equation from the previous proof reads
\[
(m_\theta-m_\tau)\sqrt{4k-3}=k^2-2k.
\]
If the square root is irrational, its integer coefficient must vanish. Then $k^2=2k$, and diameter $2$ rules out $k=0$, leaving $k=2$.
If the square root is rational, it is an integer $s>0$, necessarily odd. Substituting $k=(s^2+3)/4$ gives
\[
16(m_\theta-m_\tau)s=(s^2+3)(s^2-5)=s^4-2s^2-15.
\]
Thus $s\mid15$, so $s$ is $1,3,5$ or $15$, and $k$ is $1,3,7$ or $57$. A degree-one connected graph is a single edge, of diameter $1$. Removing that case gives the list.
\end{proof}

The pentagon supplies degree $2$, the Petersen graph degree $3$, and the Hoffman--Singleton graph degree $7$ on $50$ vertices. The degree-$57$ case would have $3250$ vertices; its existence remains open. Strong restrictions on its possible symmetries are known: Aschbacher ruled out rank-three automorphism groups, Higman ruled out vertex transitivity, and Ma\v caj and \v Sir\'a\v n bounded the automorphism group by $375$. These restrictions do not settle existence. The cited papers give their precise hypotheses and proofs.

\section{Independent sets and colourings}

An \emph{independent set} contains no adjacent pair; write $\alpha(G)$ for its largest possible size. Finding that size is hard in general. An eigenvalue bound may rule out a large set without requiring a search through all subsets.

\begin{theorem}[Hoffman's ratio bound]\label{linear:thm:hoffman}
Suppose $G$ is $k$-regular with $n$ vertices and least eigenvalue $\lambda_n<0$. Its independence number satisfies
\[
\alpha(G)\le n\frac{-\lambda_n}{k-\lambda_n}.
\]
\end{theorem}
\begin{proof}
Take an independent set $S$ of size $s$. The claim is automatic for $s=0$, so assume $s>0$. Centre its indicator by setting $x=\mathbf1_S-(s/n)\mathbf1$. Independence gives $\mathbf1_S^{\mathsf T}A\mathbf1_S=0$, and regularity gives $A\mathbf1=k\mathbf1$. Expanding both quadratic expressions yields
\[
x^{\mathsf T}Ax=-\frac{ks^2}{n},\qquad
\|x\|^2=s-\frac{s^2}{n}=\frac{s(n-s)}n.
\]
The least-eigenvalue bound holds for every vector, including zero, so
$-ks^2/n\ge\lambda_ns(n-s)/n$.
Divide by $s/n>0$ and rearrange to get $s(k-\lambda_n)\le-\lambda_nn$. The denominator $k-\lambda_n$ is positive, giving the result.
\end{proof}

For the Petersen graph the bound is $10\cdot2/(3+2)=4$. Four pairs containing a fixed element of $\{1,\ldots,5\}$ form an independent set in its disjoint-pair description, so equality occurs. That calculation extends to pairs on a larger set.

\begin{example}[Erd\H{o}s--Ko--Rado for pairs]\label{linear:exa:ekr}
Let $n\ge4$. The Kneser graph $K(n,2)$ has all $2$-element subsets of $\{1,\ldots,n\}$ as vertices, with edges joining disjoint pairs. Its independent sets are exactly the families in which any two pairs intersect.
Let $B$ be the $n\times\binom n2$ incidence matrix, whose entry $b_{i,e}$ records $i\in e$. The diagonal entries of $BB^{\mathsf T}$ are $n-1$, and its off-diagonal entries are $1$. Hence
$BB^{\mathsf T}=(n-2)I+J$, with eigenvalues $2(n-1)$ once and $n-2$ with multiplicity $n-1$. By the SVD, $B^{\mathsf T}B$ has the same non-zero eigenvalues and $\binom n2-n$ zeros.

The graph joining intersecting distinct pairs has adjacency matrix $B^{\mathsf T}B-2I$. Its complement therefore has adjacency matrix $J+I-B^{\mathsf T}B$. On $\mathbf1^\perp$ the eigenvalues of this complement are $3-n$ and $1$; its constant eigenvalue is the degree $\binom{n-2}2$. Thus Hoffman's bound gives
\[
\alpha(K(n,2))\le\binom n2
\frac{n-3}{\binom{n-2}2+n-3}
=\binom n2\frac{n-3}{n(n-3)/2}=n-1.
\]
All pairs containing one fixed element attain this bound. So a pairwise intersecting family of pairs has at most $n-1$ members. This is the pair case of the Erd\H{o}s--Ko--Rado theorem.
\end{example}

A proper colouring partitions the vertices into independent sets. Consequently $\chi(G)\ge n/\alpha(G)$, and for the regular graphs above this gives $\chi(G)\ge1+k/(-\lambda_n)$. Exercise~\ref{linear:ex:chromatic} proves a version without the regularity assumption.

\section{The Graham--Pollak theorem}

A \emph{complete bipartite graph}, or biclique, has two disjoint nonempty parts and all possible edges between them. Graham and Pollak's work on addressing telephone networks led to a question about partitioning the edges of $K_n$ into bicliques. There is a partition with $n-1$ pieces: for each $i<n$, use the star joining $i$ to all larger-numbered vertices. Linear algebra proves that no partition uses fewer pieces.

\begin{theorem}[Graham--Pollak, 1971]\label{linear:thm:grahampollak}
Every partition of the edges of $K_n$ into complete bipartite graphs has at least $n-1$ pieces.
\end{theorem}
\begin{proof}
We use Tverberg's argument. Suppose the pieces have parts $(X_t,Y_t)$ for $t=1,\ldots,m$. Since each unordered edge occurs once, every real vector $x$ satisfies
\[
\sum_{i<j}x_ix_j=\sum_{t=1}^m
\left(\sum_{i\in X_t}x_i\right)\left(\sum_{j\in Y_t}x_j\right).
\]
If $m<n-1$, the equations $\sum_{i\in X_t}x_i=0$ for all $t$, together with $\sum_ix_i=0$, form fewer than $n$ homogeneous linear equations. Rank--nullity supplies a non-zero solution. For that vector the displayed identity gives $\sum_{i<j}x_ix_j=0$, whereas
\[
0=\left(\sum_ix_i\right)^2=\sum_ix_i^2+2\sum_{i<j}x_ix_j=\sum_ix_i^2>0.
\]
This contradiction proves the lower bound.
\end{proof}

There is also a spectral reading: $J-I$ has $n-1$ negative eigenvalues, and a biclique adjacency matrix has just one. Exercise~\ref{linear:ex:inertia} proves that a sum cannot have more negative eigenvalues than the sum of those counts. Vishwanathan's 2013 paper, listed below, gives a counting proof of the same graph theorem.

\section{Huang's theorem}

The even-weight vertices of $Q_n$ form an independent set of size $2^{n-1}$; so do the odd-weight vertices. Adding a single vertex to a half-sized set must create edges, but how large a degree must it create? Chung, F\"uredi, Graham and Seymour constructed examples on this scale with maximum degree of order $\sqrt n$. Huang proved the matching lower bound.

\begin{theorem}[Huang, 2019]\label{linear:thm:huang}
Any induced subgraph of $Q_n$ with $2^{n-1}+1$ vertices contains a vertex whose degree is at least $\sqrt n$.
\end{theorem}
\begin{proof}
Put $A_1=\begin{pmatrix}0&1\\1&0\end{pmatrix}$ and recursively define
\[
A_n=\begin{pmatrix}A_{n-1}&I\\I&-A_{n-1}\end{pmatrix}.
\]
These matrices are real symmetric. Block multiplication cancels the off-diagonal blocks and gives
\[
A_n^2=\begin{pmatrix}A_{n-1}^2+I&0\\0&A_{n-1}^2+I\end{pmatrix}=nI
\]
by induction. Their eigenvalues are therefore $\pm\sqrt n$. Their traces are zero, so each sign has multiplicity $2^{n-1}$.

Order cube vertices by their last coordinate. Each diagonal block represents one copy of $Q_{n-1}$, and the identity blocks join corresponding vertices across the two copies. Thus $A_n$ has entries $\pm1$ exactly on cube edges and zeros elsewhere: it is a \emph{signed adjacency matrix}.

Let $S$ be the selected vertex set, of size $m=2^{n-1}+1$, and let $B$ be the corresponding principal submatrix. Interlacing (Theorem~\ref{linear:thm:interlacing}) gives
\[
\lambda_1(B)\ge\lambda_{1+2^n-m}(A_n)
=\lambda_{2^{n-1}}(A_n)=\sqrt n.
\]
Choose a non-zero eigenvector $y$ for $\lambda_1(B)$ and an index $i$ where $|y_i|$ is largest. In particular $|y_i|>0$. The signed entries have absolute values equal to adjacency entries of the induced graph $H$, so
\[
\lambda_1(B)|y_i|=\left|\sum_jb_{ij}y_j\right|
\le\sum_j|b_{ij}||y_j|\le\deg_H(i)|y_i|.
\]
Cancel $|y_i|$ to obtain $\deg_H(i)\ge\sqrt n$.
\end{proof}

\begin{figure}[ht]
\centering
\fig{linear-ch08-huang}
\caption{Huang's signing of the cube $Q_3$: dashed rust edges carry $-1$. Every square face has an odd number of negative edges,
which is what makes the signed matrix square to $3I$.}
\label{linear:fig:huang}
\end{figure}

\begin{remark}
For a Boolean function $f:\{0,1\}^n\to\{0,1\}$, its \emph{sensitivity} $s(f)$ is the largest number of single-coordinate flips changing its value at an input. Its \emph{block sensitivity} $bs(f)$ allows disjoint blocks of coordinates to be flipped instead. Certainly $s(f)\le bs(f)$. The sensitivity conjecture asked for an absolute exponent $C$ bounding $bs(f)$ by $s(f)^C$.
Gotsman and Linial related the problem to induced subgraphs of the cube. Combining their reduction with Huang's theorem gives $bs(f)\le2s(f)^4$. Other measures, including decision-tree depth and polynomial degree, are already polynomially related to block sensitivity. The graph theorem therefore brings sensitivity into the same group of measures. The reduction itself is not proved here; Huang's paper and the accounts by Aaronson and Tao give it.
\end{remark}

\section{A gem: three Petersens do not make a $K_{10}$}

The Petersen graph has $15$ edges and $K_{10}$ has $45$. Edge counts allow the possibility of three disjoint Petersen copies covering $K_{10}$. Their eigenspaces forbid it.

\begin{proposition}[Schwenk]\label{linear:prop:schwenk}
There is no partition of the edges of $K_{10}$ into three Petersen graphs.
\end{proposition}
\begin{proof}
Suppose their adjacency matrices satisfy $J-I=P_1+P_2+P_3$. Each Petersen matrix has a $5$-dimensional eigenspace for eigenvalue $1$, contained in the $9$-dimensional space $\mathbf1^\perp$. The eigenspaces for $P_1$ and $P_2$ must meet in a non-zero vector $x$. Then
\[
P_3x=(J-I)x-P_1x-P_2x=-3x.
\]
But the Petersen eigenvalues are $3,1,-2$, so $-3$ cannot occur. This is the desired contradiction.
\end{proof}

\begin{beauty}
The contradiction need not look like a graph. In the friendship theorem it is a divisibility condition. For Moore graphs it is the short list of divisors of $15$. In Schwenk's argument it is a vector belonging to two large subspaces. The matrix keeps the graph's rules while letting us argue with a different set of objects.
\end{beauty}

\section{How much does the second eigenvalue know?}

The party rule forced triangles around one centre. The cube argument needed a more unexpected choice: changing edge signs preserved their locations while creating a matrix with only two eigenvalues. In both cases the useful identity came from the way the matrix was built.

A regular graph always has its degree as the top adjacency eigenvalue. How far below it is the second? The next chapter relates that gap to a concrete task: separating a large set of vertices from the rest with few edges. The same gap will control how quickly a random walk approaches its stationary distribution.

\begin{figure}[ht]
\centering\fig{ch05-rw-ratio}
\caption{The four Petersen vertices corresponding to pairs containing $1$ form an independent set. The ratio bound gives exactly $4$, so no fifth pair can join this family while meeting every other pair.}
\label{fig:v8c05ratio}
\end{figure}
\begin{figure}[ht]
\centering\fig{ch05-rw-signs}
\caption{For $Q_2$, give one edge a negative sign. The two signed length-two routes between opposite vertices cancel, while each vertex has two return walks. Thus the signed adjacency matrix squares to $2I$.}
\label{fig:v8c05signs}
\end{figure}

\section{Exercises}

\begin{exercise}[$\star$]\label{linear:ex:hypercube}
For each subset $T\subseteq\{1,\ldots,n\}$, put $\chi_T(x)=(-1)^{\sum_{i\in T}x_i}$. Prove that these functions give a complete eigenbasis for $Q_n$, with eigenvalues $n-2|T|$.
\end{exercise}
\begin{exercise}[$\star$]\label{linear:ex:spectrumcount}
Suppose a graph has eigenvalues $3,1,1,1,1,1,-2,-2,-2,-2$. Determine its vertex, edge and triangle counts, and decide whether every vertex has the same degree.
\end{exercise}
\begin{exercise}[$\star\star$]\label{linear:ex:bipartite}
For a connected graph, prove the equivalence of these three conditions: it is bipartite; $-\lambda_1$ occurs in its spectrum; its spectrum is symmetric about zero, with multiplicities.
\end{exercise}
\begin{exercise}[$\star$]\label{linear:ex:srgcount}
Establish the parameter identity $k(k-\lambda-1)=(n-k-1)\mu$ for strongly regular graphs. Compute the complete spectra of the Petersen graph and of a hypothetical degree-$57$ Moore graph.
\end{exercise}
\begin{exercise}[$\star$]\label{linear:ex:moorebound}
A graph has diameter $2$ and maximum degree $k$. Prove $n\le k^2+1$. Show that equality forces regularity and excludes both triangles and four-cycles.
\end{exercise}
\begin{exercise}[$\star\star$]\label{linear:ex:linegraph}
Make a new graph $L(G)$ whose vertices are the edges of $G$, joining two when they share an endpoint. For the vertex-edge incidence matrix $B$, prove $A(L(G))=B^{\mathsf T}B-2I$. Deduce a lower bound of $-2$ for every line-graph eigenvalue.
\end{exercise}
\begin{exercise}[$\star\star$]\label{linear:ex:chromatic}
Prove Hoffman's general colouring bound: if a graph has an edge, then $\chi(G)\ge1+\lambda_1/(-\lambda_n)$. You may multiply the colour-class pieces of a Perron eigenvector by roots of unity.
\end{exercise}
\begin{exercise}[$\star\star$]\label{linear:ex:wilf}
Prove Wilf's bound $\chi(G)\le1+\lambda_1$. Use a low-degree vertex and interlacing to make an induction work on every induced subgraph.
\end{exercise}
\begin{exercise}[$\star\star\star$]\label{linear:ex:inertia}
For a biclique with nonempty parts $X,Y$, show that its adjacency matrix is $\mathbf1_X\mathbf1_Y^{\mathsf T}+\mathbf1_Y\mathbf1_X^{\mathsf T}$ and has exactly one negative eigenvalue. If $\nu(A)$ counts negative eigenvalues with multiplicity, prove $\nu(A+B)\le\nu(A)+\nu(B)$ for Hermitian matrices. Use this to recover Theorem~\ref{linear:thm:grahampollak}.
\end{exercise}
\begin{exercise}[$\star\star$]\label{linear:ex:infinitefriends}
Construct an infinite friendship graph with no vertex adjacent to everyone. Begin with three isolated vertices. At each stage, give every old pair lacking a common neighbour a new private common neighbour. Prove that the limiting graph has the required properties.
\end{exercise}
\begin{exercise}[$\star\star$]\label{linear:ex:huangsharp}
Prove that Huang's lower bound is attained for $n=1$ and $n=4$. In the latter case, find nine vertices of $Q_4$ whose induced graph has maximum degree $2$.
\end{exercise}

\begin{exercise}[$\star$]\label{ex:v8c05a}
Determine the adjacency spectrum of the complete bipartite graph $K_{a,b}$ for positive integers $a,b$, including multiplicities. Recover its edge and triangle counts from the spectrum.
\end{exercise}
\begin{exercise}[$\star\star$]\label{ex:v8c05b}
Suppose a $k$-regular graph has least eigenvalue $\lambda_n<0$, and an independent set $S$ attains Hoffman's ratio bound. Prove that $\mathbf1_S-(|S|/n)\mathbf1$ is a $\lambda_n$-eigenvector. Show that every vertex outside $S$ has exactly $-\lambda_n$ neighbours in $S$.
\end{exercise}
\begin{exercise}[$\star$]\label{ex:v8c05c}
Write the signed adjacency matrix of the square $00,10,01,11$ obtained from Huang's recursion, in that order. Verify its square by an exact multiplication and give its eigenvalues with multiplicities.
\end{exercise}
\begin{exercise}[$\star\star$]\label{ex:v8c05d}
A graph has $n$ vertices and the same spectrum as $K_n$. Prove that it is $K_n$. Then show that a graph with spectrum symmetric about zero has no odd cycle, even when it is disconnected.
\end{exercise}

\section{Hints}

\begin{hint}{linear:ex:hypercube}
Flipping coordinate $i$ changes the sign of $\chi_T$ exactly when $i\in T$. To check orthogonality, pair words differing at one coordinate of $T\triangle T'$.
\end{hint}
\begin{hint}{linear:ex:spectrumcount}
Sum the zeroth, second and third powers of the eigenvalues. Compare the constant vector's Rayleigh quotient with the largest eigenvalue.
\end{hint}
\begin{hint}{linear:ex:bipartite}
A diagonal sign change switches $A$ to $-A$ for a bipartite graph. In the converse, compare an eigenvector $y$ for $-\lambda_1$ with its entrywise absolute value.
\end{hint}
\begin{hint}{linear:ex:srgcount}
Fix a vertex $u$. Count edges from its neighbourhood to the vertices outside that neighbourhood and outside $\{u\}$.
\end{hint}
\begin{hint}{linear:ex:moorebound}
Count the three distance layers about an arbitrary vertex, keeping track of every possible repetition in the second layer.
\end{hint}
\begin{hint}{linear:ex:linegraph}
The inner product of two incidence columns counts common endpoints, including both endpoints when the columns agree.
\end{hint}
\begin{hint}{linear:ex:chromatic}
For $c$ colours and $\omega=e^{2\pi i/c}$, use $y_j=\sum_{t=0}^{c-1}\omega^{jt}x_t$ for $j=1,\ldots,c-1$. The missing $j=0$ term makes the average cross coefficient $-1/(c-1)$.
\end{hint}
\begin{hint}{linear:ex:wilf}
The average degree is a Rayleigh quotient. Delete a vertex with degree at most that average, colour what remains, and restore it.
\end{hint}
\begin{hint}{linear:ex:inertia}
Intersect the non-negative spectral subspaces of $A$ and $B$. Their intersection has codimension at most $\nu(A)+\nu(B)$.
\end{hint}
\begin{hint}{linear:ex:infinitefriends}
Check old--old, old--new and new--new pairs separately. The initial three vertices soon acquire degree greater than two.
\end{hint}
\begin{hint}{linear:ex:huangsharp}
Find an induced eight-cycle in the four-cube and a ninth vertex with no neighbour on that cycle.
\end{hint}
\begin{hint}{ex:v8c05a}
The adjacency map preserves the two-dimensional space of vectors constant on each part and kills its orthogonal complement.
\end{hint}
\begin{hint}{ex:v8c05b}
Equality in the least-eigenvalue bound puts the centred indicator in the least eigenspace. Read its eigenvector equation at a vertex outside the set.
\end{hint}
\begin{hint}{ex:v8c05c}
Use $A_2=\begin{pmatrix}A_1&I_2\\I_2&-A_1\end{pmatrix}$ rather than guessing the signs independently.
\end{hint}
\begin{hint}{ex:v8c05d}
For the first claim, count edges. For the second, an odd cycle gives a positive contribution to an odd trace, whereas a symmetric spectrum makes every odd trace vanish.
\end{hint}

\section{Solutions}

\begin{solution}{linear:ex:hypercube}
\idea Count the signs contributed by the $n$ coordinate flips.
\writeup Neighbours of $x$ are $x+e_i$, with addition modulo $2$. Of the terms in $(A\chi_T)(x)=\sum_i\chi_T(x+e_i)$, exactly $|T|$ equal $-\chi_T(x)$ and the other $n-|T|$ equal $\chi_T(x)$. The sum is $(n-2|T|)\chi_T(x)$.
If $T\ne T'$, choose $i\in T\triangle T'$. Pairing each word with its $i$th-coordinate flip cancels the sum of $\chi_T\chi_{T'}$. Thus the $2^n$ non-zero vectors are mutually orthogonal and form a basis. There are $\binom nj$ subsets of size $j$, proving the multiplicity formula as well.
\end{solution}

\begin{solution}{linear:ex:spectrumcount}
\idea The traces give counts, and equality in a Rayleigh bound gives regularity.
\writeup Ten eigenvalues mean ten vertices. Their squares sum to $9+5+16=30$, so there are $15$ edges. Their cubes sum to $27+5-32=0$, so there are no triangles. The average degree is $30/10=3$. It equals $\lambda_1$, so the constant vector attains the maximum Rayleigh quotient. Expanding it in an orthonormal eigenbasis shows it belongs to the $3$-eigenspace, as in Exercise~\ref{linear:ex:degree}. Therefore $A\mathbf1=3\mathbf1$ and every degree is $3$.
\end{solution}

\begin{solution}{linear:ex:bipartite}
\idea A bipartition changes edge signs globally; an extreme negative eigenvector recovers that partition.
\writeup For a single vertex, all three conditions hold. Otherwise connectedness gives an edge and $\lambda_1>0$. If the graph is bipartite, let $D$ have diagonal entries $1$ on one side and $-1$ on the other. Then $DAD=-A$. Similar matrices have the same eigenvalues with multiplicities, proving symmetry about zero and in particular the occurrence of $-\lambda_1$.
Conversely take a real non-zero $y$ with $Ay=-\lambda_1y$. Entrywise non-negativity of $A$ gives
\[
\lambda_1\|y\|^2=-y^{\mathsf T}Ay
\le |y|^{\mathsf T}A|y|\le\lambda_1\|y\|^2.
\]
Equality in the last Rayleigh bound makes $|y|$ a top eigenvector. Frobenius's theorem (Volume~3, Theorem~A.6) makes it strictly positive in a connected graph. Equality in the first inequality forces $y_iy_j<0$ on every edge: each non-negative difference term must vanish, and no coordinate is zero. The positive and negative coordinates therefore form a bipartition. Finally, spectral symmetry always implies the occurrence of $-\lambda_1$, completing all implications.
\end{solution}

\begin{solution}{linear:ex:srgcount}
\idea Count edges between two distance layers, then solve the quadratic from the matrix identity.
\writeup Fix $u$. Each of its $k$ neighbours has $k-1-\lambda$ neighbours outside $\{u\}\cup N(u)$; each of the $n-k-1$ vertices there has $\mu$ neighbours in $N(u)$. Counting the same edges gives $k(k-1-\lambda)=(n-k-1)\mu$.
For Petersen, the quadratic is $x^2+x-2$, with roots $1,-2$. The equations $m_1+m_{-2}=9$ and $3+m_1-2m_{-2}=0$ give multiplicities $5,4$.
For degree $57$, $n=57^2+1=3250$ and $\sqrt{4k-3}=15$. The roots are $7,-8$. Solving $m_7+m_{-8}=3249$ and $57+7m_7-8m_{-8}=0$ gives $m_7=1729$ and $m_{-8}=1520$. These integral multiplicities are necessary conditions, not an existence proof.
\end{solution}

\begin{solution}{linear:ex:moorebound}
\idea Equality must hold in every layer count, about every choice of centre.
\writeup From any vertex $u$, there are at most $k$ first neighbours and at most $k(k-1)$ further vertices reached through them. Diameter $2$ exhausts the graph in these layers, so $n\le1+k+k(k-1)=k^2+1$. If equality holds, the first layer has exactly $k$ vertices, each contributes $k-1$ new vertices, and all those contributions are distinct and outside the first layer. Thus $\deg u=k$. A triangle through $u$ would send a contribution back into the first layer, and a four-cycle through $u$ would repeat a second-layer vertex. Both are excluded. Since $u$ was arbitrary, the whole graph is regular and has neither type of cycle.
\end{solution}

\begin{solution}{linear:ex:linegraph}
\idea A Gram matrix is positive semidefinite, regardless of the graph used to build its columns.
\writeup The column for edge $e=\{i,j\}$ has ones precisely at $i,j$. Thus $(B^{\mathsf T}B)_{ef}=|e\cap f|$. This is $2$ on the diagonal and the line-graph adjacency entry elsewhere, proving the formula. If $z$ is an eigenvector of $A(L(G))$ with eigenvalue $\lambda$, then
$(\lambda+2)\|z\|^2=z^{\mathsf T}B^{\mathsf T}Bz=\|Bz\|^2\ge0$.
Since $z\ne0$, $\lambda\ge-2$.
\end{solution}

\begin{solution}{linear:ex:chromatic}
\idea Average several phase choices so every cross-class term has the same coefficient.
\writeup Let $c=\chi(G)\ge2$ and label the colour classes $0,\ldots,c-1$. Choose a unit non-negative top eigenvector $x$; this exists by the Perron theory in Volume~3, also for a disconnected graph by choosing a component with largest eigenvalue. Write $x_t$ for its restriction to class $t$, extended by zero. The pieces are orthogonal, their squared norms sum to $1$, and $x_t^{\mathsf T}Ax_t=0$.
Put $\omega=e^{2\pi i/c}$ and $y_j=\sum_t\omega^{jt}x_t$. Each $y_j$ has norm $1$, so $y_j^*Ay_j\ge\lambda_n$. For $t\ne t'$, the geometric sum gives $\sum_{j=1}^{c-1}\omega^{j(t'-t)}=-1$. Averaging the Rayleigh numerators over these $j$ therefore gives
\[
\lambda_n\le-\frac1{c-1}\sum_{t\ne t'}x_t^{\mathsf T}Ax_{t'}
=-\frac{\lambda_1}{c-1}.
\]
A graph with an edge has $\lambda_1>0$ and, since its trace is zero, $\lambda_n<0$. Rearrangement gives $c\ge1+\lambda_1/(-\lambda_n)$.
\end{solution}

\begin{solution}{linear:ex:wilf}
\idea The Rayleigh bound supplies a vertex that fits into a palette of $\lfloor\lambda_1\rfloor+1$ colours.
\writeup Induct on the number of vertices, the one-vertex case being immediate. Some vertex $v$ has degree no larger than the average degree, and
$2|E|/n=R_A(\mathbf1)\le\lambda_1(G)$.
Thus $\deg v\le\lfloor\lambda_1(G)\rfloor$. Interlacing gives $\lambda_1(G-v)\le\lambda_1(G)$. By induction, $G-v$ can be coloured with at most $\lfloor\lambda_1(G)\rfloor+1$ colours. At most $\lfloor\lambda_1(G)\rfloor$ are used on neighbours of $v$, so one remains for $v$. This proves the stronger integer-valued bound and hence the stated inequality.
\end{solution}

\begin{solution}{linear:ex:inertia}
\idea A negative spectral subspace cannot meet a subspace where the quadratic form is non-negative.
\writeup Put $u=\mathbf1_X$, $v=\mathbf1_Y$. These are non-zero and orthogonal. The biclique matrix $uv^{\mathsf T}+vu^{\mathsf T}$ has eigenvectors $u/\|u\|\pm v/\|v\|$ with eigenvalues $\pm\|u\|\|v\|$ and vanishes on their orthogonal complement. It has exactly one negative eigenvalue.
For subadditivity, let $N_A,N_B$ be the spans of all non-negative eigenvectors of $A,B$. Their codimensions are $a=\nu(A)$ and $b=\nu(B)$, so $\dim(N_A\cap N_B)\ge n-a-b$. Both forms are non-negative on that intersection, and hence so is their sum. If the negative eigenspace of $A+B$ had dimension greater than $a+b$, it would meet the intersection nontrivially, contradicting strict negativity there. This proves the bound; if $a+b\ge n$, it is already automatic. Summing biclique matrices to $J-I$ now gives $n-1=\nu(J-I)\le m$.
\end{solution}

\begin{solution}{linear:ex:infinitefriends}
\idea Repair missing common neighbours without ever repairing the same old pair twice.
\writeup Let $G_0$ consist of isolated vertices $a,b,c$. Given the finite graph $G_k$, add for every unordered pair $\{u,v\}$ with no common neighbour a distinct vertex $w_{uv}$ joined only to its two parents $u,v$. Do not join new vertices to each other. Inductively assume that no old pair has two common neighbours.
An old pair gets a new common neighbour exactly when it lacked one, and then gets only its own new vertex. For an old vertex $x$ and a new $w_{uv}$, common neighbours must be parents $u,v$; both cannot qualify, since then $x$ would already have been a common neighbour of the parent pair. Two distinct new vertices have at most one parent in common, so they too have at most one common neighbour. This proves the induction.
In the union $G=\bigcup_kG_k$, any two vertices appear together at a finite stage and have a common neighbour by the next stage. They cannot acquire two, since those vertices would all occur together at some finite stage. Thus every pair has exactly one.
To rule out a universal vertex, observe that $G_1$ is the hexagon $a,w_{ab},b,w_{bc},c,w_{ca}$. At the next stage, each of the pairs $\{a,w_{ab}\}$, $\{a,w_{ca}\}$ and $\{a,w_{bc}\}$ receives a new neighbour. So $a$ has degree at least $5$, and the same holds for $b$. If a universal vertex existed, the last paragraph of the finite friendship proof would force every other vertex to have degree exactly $2$. At least one of $a,b$ is different from that vertex, a contradiction. This also shows that the union cannot be finite, since the finite theorem would then supply a universal vertex.
\end{solution}

\begin{solution}{linear:ex:huangsharp}
\idea Use a chordless eight-cycle and one isolated vertex to reach nine vertices with maximum degree two.
\writeup For $n=1$ the required two vertices form the single cube edge, of degree $1$. For $n=4$, take in cyclic order
\[
0000,0001,0101,0111,1111,1011,1010,0010
\]
and add $1100$. To verify the induced edges without a search, list each cycle vertex's neighbours lying in this set:
\[
\begin{array}{c|c@{\qquad}c|c}
0000&0001,0010&0001&0000,0101\\
0101&0001,0111&0111&0101,1111\\
1111&0111,1011&1011&1111,1010\\
1010&1011,0010&0010&1010,0000
\end{array}
\]
Each list follows by flipping its four bits. The neighbours of $1100$ are $0100,1000,1101,1110$, all absent. These exact lists prove that the induced graph is an eight-cycle plus an isolated vertex, attaining degree $2=\sqrt4$.
\end{solution}

\begin{solution}{ex:v8c05a}
\idea Separate the two constant directions from vectors summing to zero on each part.
\writeup Let $u=\mathbf1_X/\sqrt a$ and $v=\mathbf1_Y/\sqrt b$. They are orthonormal, and adjacency sends $u$ to $\sqrt{ab}\,v$ and $v$ to $\sqrt{ab}\,u$. Thus $u+v$ and $u-v$ have eigenvalues $\sqrt{ab}$ and $-\sqrt{ab}$. A vector orthogonal to both has zero coordinate sum in each part, so adjacency kills it. The spectrum is therefore these two eigenvalues and $0$ with multiplicity $a+b-2$. Its squared sum is $2ab$, giving $ab$ edges; its cubed sum is zero, giving no triangles.
\end{solution}

\begin{solution}{ex:v8c05b}
\idea Equality in a weighted eigenvalue average excludes every larger eigenspace.
\writeup Write $s=|S|$. Since $k>0$ and $\lambda_n<0$, equality in the bound gives $0<s<n$. The non-zero vector $x=\mathbf1_S-(s/n)\mathbf1$ attains $x^{\mathsf T}Ax=\lambda_n\|x\|^2$ in the proof. Expanding $x$ in an orthonormal eigenbasis makes this identity
$\sum_i(\lambda_i-\lambda_n)|c_i|^2=0$.
Every summand is non-negative, so only the least-eigenvalue coefficients remain. Hence $Ax=\lambda_nx$.
At a vertex $v\notin S$, this equation says
$|N(v)\cap S|-ks/n=-\lambda_ns/n$.
The equality assumption also gives $s(k-\lambda_n)/n=-\lambda_n$, so $|N(v)\cap S|=-\lambda_n$, independently of $v$.
\end{solution}

\begin{solution}{ex:v8c05c}
\idea Keep the recursive block order to make the cancellation visible.
\writeup The requested matrix is
\[
A_2=\begin{pmatrix}0&1&1&0\\1&0&0&1\\1&0&0&-1\\0&1&-1&0\end{pmatrix}.
\]
Each row has squared length $2$. Distinct rows have inner product zero: rows with disjoint support contribute zero, and the overlapping pairs contribute $1-1$. Thus $A_2A_2^{\mathsf T}=2I_4$; symmetry gives $A_2^2=2I_4$. All eigenvalues are $\pm\sqrt2$, and the zero trace makes their multiplicities two each.
\end{solution}

\begin{solution}{ex:v8c05d}
\idea Closed-walk counts detect missing edges and forbid odd cycles under spectral symmetry.
\writeup The spectrum of $K_n$ is $n-1,-1,\ldots,-1$, so its squared sum is $(n-1)^2+(n-1)=n(n-1)$. A graph with that spectrum has $n(n-1)/2$ edges, the total number of unordered pairs of distinct vertices. Every pair must therefore be an edge.
If a spectrum is symmetric, then $\operatorname{tr}A^\ell=0$ for every odd positive integer $\ell$, by cancellation of opposite eigenvalues. An odd cycle of length $\ell$ would give at least $2\ell$ closed walks of that length, so $\operatorname{tr}A^\ell>0$, since all walk counts are non-negative. This contradiction excludes odd cycles in every component.
\end{solution}

\section*{Sources and further reading}

\begin{itemize}
\item P.~Erd\H{o}s, A.~R\'enyi and V.~T.~S\'os, \emph{On a problem of graph theory}, Studia Scientiarum Mathematicarum Hungarica 1
(1966), 215--235; M.~Aigner and G.~M.~Ziegler, \emph{Proofs from THE BOOK}, 6th ed., Springer, 2018, chapters ``Of friends and
politicians'' and ``Three famous theorems on finite sets'' (Graham--Pollak).
\item A.~J.~Hoffman and R.~R.~Singleton, \emph{On Moore graphs with diameters 2 and 3}, IBM Journal of Research and Development 4 (1960),
497--504; M.~Aschbacher, \emph{The non-existence of rank three permutation groups of degree 3250 and subdegree 57}, Journal of Algebra 19
(1971), 538--540; M.~Ma\v caj and J.~\v Sir\'a\v n, \emph{Search for properties of the missing Moore graph}, Linear Algebra and its
Applications 432 (2010), 2381--2398, which also gives Higman's argument against vertex-transitivity.
\item A.~J.~Hoffman, \emph{On eigenvalues and colorings of graphs}, in \emph{Graph Theory and its Applications}, Academic Press, 1970,
79--91; H.~S.~Wilf, \emph{The eigenvalues of a graph and its chromatic number}, Journal of the London Mathematical Society 42 (1967),
330--332.
\item R.~L.~Graham and H.~O.~Pollak, \emph{On the addressing problem for loop switching}, Bell System Technical Journal 50 (1971),
2495--2519; H.~Tverberg, \emph{On the decomposition of $K_n$ into complete bipartite graphs}, Journal of Graph Theory 6 (1982), 493--494; S.~Vishwanathan, \emph{A counting proof of the Graham--Pollak theorem}, Discrete
Mathematics 313 (2013), 765--766.
\item A.~J.~Schwenk, \emph{Problem 6434}, American Mathematical Monthly 90 (1983), 403, and solution 92 (1985), 591--592 (three Petersens).
\item H.~Huang, \emph{Induced subgraphs of hypercubes and a proof of the Sensitivity Conjecture}, Annals of Mathematics 190 (2019),
949--955; F.~R.~K.~Chung, Z.~F\"uredi, R.~L.~Graham and P.~Seymour, \emph{On induced subgraphs of the cube}, Journal of Combinatorial
Theory A 49 (1988), 180--187; C.~Gotsman and N.~Linial, \emph{The equivalence of two problems on the cube}, Journal of Combinatorial Theory
A 61 (1992), 142--146; N.~Nisan and M.~Szegedy, \emph{On the degree of Boolean functions as real polynomials}, Computational Complexity 4
(1994), 301--313.
\item S.~Aaronson, \emph{Sensitivity conjecture resolved}, blog post, 2 July 2019; T.~Tao, \emph{Twisted convolution and the sensitivity
conjecture}, blog post, 2 July 2019.
\item A.~E.~Brouwer and W.~H.~Haemers, \emph{Spectra of Graphs}, Springer, 2012; C.~Godsil and G.~Royle, \emph{Algebraic Graph Theory},
Springer, 2001, Chapters 5, 9 and 10.
\end{itemize}
